\begin{theorem}[Soundness and completeness of linear temporal logic]
For every temporal formula \(A\),
\[
  \derives_{\mathsf{LTL}}A
  \quad\Longleftrightarrow\quad
  \mathcal M,n\entails A
  \text{ for every temporal structure \(\mathcal M\) and every \(n\in\N\)}.
\]
This is weak completeness for validity, not a strong-completeness assertion
for arbitrary premise sets.  It concerns discrete infinite linear time;
finite traces, dense time, and branching time require different semantics and
calculi.
\end{theorem}

\paragraph{Complete proof.}
The seven registered Hilbert rules, including theorem-only next necessitation
and until induction, are sound by Theorem~\ref{thm:p08-ltl-soundness}.
Theorem~\ref{thm:p08-linear-temporal-completeness} and
Annex~\ref{annex:p08.ltl.completeness} supply consistent finite types,
successor existence, realization of every pending until obligation and the
complete infinite-word truth lemma. Until here admits the current position
as its witness. The strict-future language and source calculus are compared
separately; their symbols and axioms are not silently substituted. The result
is weak completeness for validity, with finite-premise consequence treated
separately. It is not strong completeness for arbitrary premise sets.
